\section{Conclusion}\label{conclusion}

Observed hours and earnings do not say whether a household chose its position or settled for it, and that ambiguity is not a nuisance for welfare measurement: it is the substance of it. This paper takes the ambiguity seriously in both halves of the problem. On the behavioural side it estimates a model in which the jobs a household can reach and the way it ranks them are identified jointly, with every alternative priced through the tax-benefit system. On the normative side it evaluates well-being with a reference that keeps each household's own preferences and own reachable jobs and removes variation in pay from the reference bundles.

The decomposition then asks how the resulting inequality divides. For French single-adult households, labour-market opportunities carry 55.92 per cent of it and job access alone 49.16 per cent, more than household resources and needs and far more than preferences. For couples the same accounting gives a different answer: earning opportunities and household circumstances dominate and access matters little. The ordering of access against earnings within each type is the part that survives all six inequality indices we report; the preference contribution is not, and changes sign outside the Gini.

Three limits should travel with those numbers. The measure is one member of a family of references and the shares move with the reference, in ways we report rather than resolve. The allocation is exhaustive for the declared game, which validates the accounting and not the identification or the ethics of the operators. And two questions about the observation rule remain genuinely open. What the exercise offers is not a causal account of why opportunities differ, but a disciplined statement of how much of measured well-being inequality is associated with them once preferences, resources and needs are modelled explicitly and the reference is stated.

